%% Kaynak: tez-latex-v2/frontmatter/ozet.tex (15.11.2026)
\begin{ozet}
Enerji depolama sistemlerinin başarısı yalnız batarya kapasitesiyle değil, güç dönüşüm arayüzünün verimi, güvenilirliği ve maliyetiyle de belirlenir. Şebekeye bağlı üç seviyeli T-tipi evirici, anma gücünde yüksek verim ve güç yoğunluğu sunmasına karşılık çalışma süresinin önemli bir bölümünü oluşturan düşük güç bölgesinde farklı bir sorunla karşılaşır: sıfır gerilim durumunu üreten nötr anahtar kolu, endüktör akımı sıfıra ulaştıktan sonra iletimde tutulduğunda senkron alçaltıcı dönüştürücü benzeri bir ters akım dolaşımı oluşturur. Yararlı güç taşımayan dolaşım, RMS akımı ve iletim kayıplarını artırarak hafif yük verimini düşürür.

Tezde ters akım mekanizması anahtarlama durumları, akım yolları ve senkron-buck eşdeğeri üzerinden çözümlenmiş ve T-tipi topolojiye dayanan üç özgün katkı tek bir araştırma zinciri altında birleştirilmiştir. Birinci katkı, ters akımı sürdüren nötr anahtar komutunu düşük güç bölgesinde bastıran nötr anahtar susturma yöntemidir. Yaklaşık 2,3~A faz akımında yöntem, klasik anahtarlamaya göre verimi deneyde yüzde 92,5'ten yüzde 94,2'ye, deneyin aktif ve reaktif gücüne kalibre edilmiş benzetimde yüzde 91,15'ten yüzde 93,93'e yükseltmiştir; karşılığında akımı süreksiz iletime zorlayarak şebeke akımı THD'sini deneyde yüzde 15'ten yüzde 47'ye çıkarmıştır. İkinci katkı, harmonik cezayı azaltmak için şebeke faz bilgisinden üretilen Sin\textsuperscript{3} tabanlı ileri besleme kompanzasyonudur. Susturma ile birlikte kullanılan terim deneyde THD'yi yüzde 14'e, yani klasik anahtarlama düzeyine indirmiş ve verimi yüzde 94,5 düzeyinde korumuştur. Kalibre edilmiş benzetim modeli, kalibrasyon hedefi olmayan THD'yi yüzde 14,8, verimi 0,25 puan içinde yeniden üretmiştir. $\sin^3\theta=[3\sin\theta-\sin(3\theta)]/4$ özdeşliğine dayanan parçalı DCM--CCM modeliyle yapılan tarama, katsayıların çalışma noktasına bağlı olduğunu göstermiş ve $I_{\mathrm{rms}}>\SI{2.2}{\ampere}$ koşulunda $(k_1,k_3)=(0{,}36,\ 0{,}42)$ çiftiyle üçüncü harmonik oranını analitik modelde yüzde 39,96'dan yüzde 0,315'e indirmiştir. Akım geri beslemeli kapalı çevrim benzetiminde PI denetimi 25~A anma akımını eksi yüzde 2,1 genlik hatasıyla, doyumsuz izlemiştir. Kararlılık analizi, Sin\textsuperscript{3} teriminin karakteristik denklemde yer almadığını ve çevrim kazancı ile duyarlılık fonksiyonlarını değiştirmediğini göstermiştir. Analog eşdeğer olarak modellenen akım çevrimi 54° faz marjı ve 35~dB kazanç marjıyla kararlıdır. Üçüncü katkı, aynı T-tipi yarı iletkenlerin evirici, Vienna PFC ve aktif batarya yönetimi işlevlerinde yeniden görevlendirilmesine dayanan GT-tipi yeniden yapılandırılabilir enerji depolama mimarisidir. Deney düzeneği GT-tipi olarak kurulmuştur. Düşük güç deneyleri donanımın evirici kipinde yapılmış, şarj yönünün işlevsel bütünlüğü benzetim kayıtlarıyla gösterilmiş, çalışma kipleri ve güvenli kip geçiş koşulları tanımlanmıştır. Mimari, üç işlevi ortak yarı iletkenlerde birleştirerek eleman sayısını azaltma olanağı sunar.

Benzetim modelleri tez tanımına göre denetlenmiş, tanıma uyarlanmış ve deney özetine kalibre edilmiştir. Benzetim ve deney sonuçları ayrı raporlanmıştır. Deneyler, güvenlik amacıyla varyak üzerinden yaklaşık 120~V şebeke geriliminde yürütülmüştür. Orta ve anma akımı sonuçları kapalı çevrim benzetimiyle elde edilmiştir.

\textbf{Anahtar Kelimeler:} T-tipi evirici, enerji depolama sistemi, düşük güç, nötr anahtar susturma, ileri besleme kompanzasyonu, üçüncü harmonik, yeniden yapılandırılabilir mimari.
\end{ozet}
